% !TEX program = xelatex
%---------------------------------------------------------------------------
% Paper: Transmission dynamics of the 2026 Bundibugyo ebolavirus disease
%        outbreak in the Democratic Republic of the Congo
% Status: revision draft with v6 numeric values (calibration window extended
%         through 20 Oct 2026 on the verified UKHSA weekly series wk21-wk41
%         plus the WHO DON603 declaration anchor); referee-report revisions
%         of 2026-10-23 applied.
% Target: PLOS Computational Biology (document class to be swapped to the
%         PLOS template at submission; Author Summary / statements below are
%         template-independent).
% Submission checklist (complete before submission):
%   [ ] swap documentclass to PLOS template
%   [ ] insert repository URL + Zenodo DOI in Code and data availability
%   [x] insert report numbers / URLs for surveillance citations
%   [x] confirm Funding statement wording (four PLOS statements in frontmatter)
%   [ ] attach Supplementary Material (verified weekly series + scripts)
% Language: English (discussion with the team is in Chinese).
%---------------------------------------------------------------------------
\documentclass[10pt]{plos}
\usepackage{amsmath,amssymb,amsthm}
\usepackage{xcolor}
\usepackage{booktabs}
\usepackage{tabularx}
% \usepackage{xltabular}  % unused in body; removed for PLOS build
\usepackage{graphicx}
\graphicspath{{./}}
\usepackage{url}
\usepackage{xurl}
\usepackage[colorlinks=true,linkcolor=blue,citecolor=blue]{hyperref}
\usepackage{caption}

\newtheorem{proposition}{Proposition}
\newcommand{\Reff}{\mathcal{R}}
% Zenodo archived-snapshot DOI -- replace XXXXXXX with the real deposited DOI
\newcommand{\zenododoi}{10.5281/zenodo.XXXXXXX}

\begin{document}
\title{Transmission dynamics of the 2026 Bundibugyo ebolavirus disease
outbreak in the Democratic Republic of the Congo:\\
a compartmental model with an explicit observation layer}

\author{Tang Shujiang}
\affil{School of Science, Hunan Institute of Science and Engineering, Yongzhou, China.}
\email{sjtang@huse.edu.cn}
\maketitle

\begin{abstract}
A Bundibugyo ebolavirus outbreak that began in eastern Democratic
Republic of the Congo in May 2026 has become the largest epidemic ever
caused by this species. Confirmed cases initially grew far faster
($0.109\,\mathrm{d}^{-1}$) than the virus's natural history can sustain,
so the reported curve conflates transmission with detection. We develop a
path-mixed SEIHRD model---symptomatic individuals split at onset into
care-seeking and community tracks---coupled to an observation layer in
which a logistic detection ramp maps infections to confirmed cases.
Calibrated on a verified 21-week UKHSA series (wk21--wk41) anchored at
the outbreak declaration and closed by four daily register intervals
under a negative-binomial likelihood, it reconciles the early rise with
slower intrinsic growth ($\Reff_0\approx1.73$, 95\% interval
$[1.44,1.98]$; current $\Reff$ below one from about day 115). Constant
detection worsens the AIC by 42 and overstates the inferred
transmissibility by 54\%, while the detection ceiling carries no
information from case counts and must be fixed from independent anchors.
Against Uganda 2007, the only other Bundibugyo event with a quantified
reproduction number ($\Reff_t=1.55$), the present transmissibility lies
inside that outbreak's credible interval: the larger burden reflects
context, not a more transmissible virus. Under uniform mixing, ring
vaccination's effect on cumulative deaths hinges on the unobserved
fraction of transmission inside traceable rings.
\end{abstract}

\paragraph*{Keywords}
Bundibugyo ebolavirus; compartmental model; observation model; effective
reproduction number; negative binomial likelihood; ring vaccination;
Democratic Republic of the Congo.

% (Highlights: Elsevier-only element removed for PLOS; see Author Summary.)

\paragraph*{Author summary}
Confirmed-case counts conflate how many people are infected with how
many infections are detected. We built a model that separates the two
for the 2026 Bundibugyo outbreak in the Democratic Republic of the
Congo. The virus is not unusually transmissible---its reproduction
number, about 1.73, sits inside the credible interval of the 2007
Uganda outbreak---but delayed detection, conflict and a high share of
community deaths let the outbreak grow far larger. Current transmission
is estimated below one, yet the projection band spans a factor of 2.2,
so the response is best judged on raising detection and cutting
community transmission, above all through safe and dignified burial.
\paragraph*{Ethics statement}
This study is a secondary analysis of publicly available, aggregated
surveillance data. No individual-level data were collected, no human
participants were involved, and no interventions were performed; the
analysis therefore did not require research ethics committee approval.

\paragraph*{Funding statement}
The author received no specific funding for this work.

\paragraph*{Competing interests}
The author has declared that no competing interests exist.

\paragraph*{Author contributions}
Conceptualization, data curation, formal analysis, investigation,
methodology, software, validation, visualization, writing---original
draft, and writing---review and editing: Tang Shujiang. The author is the
sole author and performed all roles.

\paragraph*{Data availability}
All inputs are public surveillance series published by WHO, UKHSA, ECDC, Africa CDC and the DRC Ministry of Public Health (enumerated with URLs in the deposited repository and Supplementary Material);
the model code, calibration and scenario pipeline, the verified weekly
series and all table/figure scripts are deposited in a public
version-controlled repository under an open license (URL and
archived-snapshot DOI \zenododoi{} in section~\ref{sec:availability}); the weekly
series and data-processing scripts are also Supplementary Material.

\section{Introduction}\label{sec:intro}

Ebola virus disease is caused by six recognized \emph{Orthoebolavirus}
species, of which the Zaire species has driven every large outbreak on
record---the 2014--2016 West African epidemic and the 2018--2020 North
Kivu--Ituri epidemic among them---and consequently attracts almost all
quantitative transmission modeling; it also established the decisive
role of case isolation and safe burial \cite{Pandey2014,Kucharski2015,Merler2015,Meltzer2014}.
The Bundibugyo species, first identified in Uganda in 2007
\cite{Wamala2010}, had by 2026 caused only two recognized human
epidemics of any size (Uganda 2007: 149 cases, 37 deaths; DRC 2012: 57
cases, 29 deaths), and its catalogue of transmission estimates is,
outside a single reconstruction of the 2007 event \cite{DePadua2026},
essentially empty. None of this provides a calibrated transmission
model for Bundibugyo virus in a contemporary setting.

That gap matters because the 2026 Bundibugyo outbreak in the eastern DRC
has, by cumulative case count, become the largest epidemic ever caused
by this species: declared in Ituri Province on 15 May 2026, it had reached 9{,}415 confirmed cases and
4{,}536 deaths across 66 health zones in eight provinces by 16 October. Three features make it distinctive: it unfolds in an
active conflict zone with large-scale displacement \cite{SSHAP2026}; a
large share of deaths occurs in the community (71\%
in the period ending 8 October); and an
occupational vaccination campaign using the Zaire-species vaccine Ervebo
is under way, with 7{,}775 frontline workers vaccinated under a
compassionate-use protocol and 848 more under the (subsequently
suspended) BRAVO protocol by 7 October.

We address the gap with a path-mixed SEIHRD model coupled to an explicit
observation layer. Three findings follow. First, the detection
ramp---not intrinsic transmission---absorbs the early confirmed-case
acceleration: constant detection worsens the AIC by 42 and overstates
the inferred $\Reff_0$ by 54\%. Second, the growth-phase $\Reff_0=1.73$
lies inside the 2007 Uganda estimate's credible interval ($1.55$, 95\%
CrI $1.02$--$2.36$) and at the lower end of the Zaire range, so the
outbreak's scale reflects context rather than a more transmissible
virus. Third, the case-count likelihood carries no information about the
detection ceiling, one-year projections are intervals spanning a factor
of 2.2, and the ring-vaccination effect on cumulative deaths is small
and hinges on the unobserved fraction of transmission inside traceable
rings.
\section{Data and methods}\label{sec:methods}

\subsection{Data sources and limitations}\label{sec:data}

All quantitative statements draw on a single verified source family for
the calibration series: the UK Health Security Agency \emph{Outbreaks
under monitoring} reports wk21--wk41 (Table~\ref{tab:weekly}),
transcribed from the national situation reports of the Congolese
public-health institute. Every weekly increment reproduces the
cumulative difference of the preceding report exactly (15/15 adjacent
pairs from wk21 to wk36, and the wk36$\to$wk37 increment against the
published ``increase of 501 confirmed cases and 243 deaths''), adjacent
data dates are 6--8 days apart, and the computed case fatality agrees
with the source-stated value to within 0.5 percentage points at all 15
checkable dates. The series opens at the declaration anchor (83
confirmed cases and 9 deaths on 21 May) and runs to
9{,}066/4{,}381 on 11 October; four daily register intervals close the
calibration window (9{,}277/4{,}478 on 14 October;
9{,}415/4{,}536 on 16 October; and
9{,}668/4{,}649 on 20 October; the WHO register for the same date lists 9{,}704/4{,}687). The DRC national
register runs about 88 cases above the UKHSA register at the
11-October splice and is not spliced with it.

Mid-outbreak WHO anchors interleave with the weekly points and serve as
independent cross-checks rather than likelihood inputs: the
Director-General's briefing of 3 June (344/60); the
Disease Outbreak News of 10 June, 17 June, 1 July, 15 July, 30 July, 12
August and 26 August (676/136, 896/232, 1{,}460/452, 2{,}124/828,
3{,}605/1{,}587, 4{,}665/2{,}184 and 5{,}794/2{,}786, respectively
);
and the mid-August snapshot of 5{,}290/2{,}516, whose data date is 19
August in the ECDC report and 21 August in the WHO commentary. Each falls on or between weekly points and
agrees with the UKHSA backbone to within the documented cross-source
noise. They are kept out of the likelihood because adjacent
cross-source snapshots disagree in both directions---363 confirmed on 2
June against 344 on 3 June, and 7{,}773/3{,}759 against 7{,}733/3{,}732
on 20 September---so no source family can be
extended by borrowing another family's points.

Two one-off registration transfers contaminate the increment series and
are handled explicitly. First, the week ending 27 July contains a
retrospective reconciliation: of the 887 new confirmed cases and 488 new
deaths booked between 20 and 27 July, 272 cases (30.7\%) and 174 deaths
(35.7\%) were ``identified through retrospective data reconciliation and
may not represent new infections or deaths that occurred during the
current reporting period''---a 22 July backfill
corroborated by the WHO figure note  and by the African
Regional Office report of the following week, which attributes its
smaller increase to ``the absence of the large retrospective data
reconciliation that contributed to last week's increase''. Because the onset dates of the reconciled cases are
published by no source, the three affected weekly intervals are merged
into one 21-day interval (13 July--3 August; 1{,}863 cases, 88.7 per
day), diluting the contamination from 30.7\% to 14.6\% of the interval
without any numerical correction: the three constituent weekly increments
(mid-points $t=53$, $60$, $67\to74$) are collapsed into a single
21-day block weighted equally by day (1{,}863/21 = 88.7 per day), so the
reconciled share is diluted rather than imputed to onset dates. The raw
increment and the
reconciliation-excluded increment are reported as sensitivities
(section~\ref{sec:ident}). Second, the first two intervals inherit a
suspected-to-confirmed reclassification: the suspected count fell from
1{,}077 (26 May) to 906 (27 May) and from ``over 1000'' to 116 in the
first week of June as the testing backlog was cleared, so part of the early
apparent growth is register transfer rather than new infection---
precisely the signal the observation layer of
section~\ref{sec:obsmodel} absorbs. Merging the first twelve days into a
single interval (280 cases) or dropping them moves $\Reff_0$ by $+1.4\%$
or $-1.4\%$ respectively (section~\ref{sec:ident}).
\begin{table}[htbp]
\centering
\caption{Verified weekly series: UKHSA \emph{Outbreaks under monitoring}
reports wk21--wk41 at their DRC data dates, with the 21 May declaration
anchor. Values are confirmed cases unless noted; weekly increments are
differences of consecutive data dates (6--8 days apart).}
\label{tab:weekly}
\footnotesize
\begin{tabular}{lrrrrr}
\toprule
Week ending & $t$ (d) & Cum.\ confirmed & Incr. & Cum.\ deaths & Death incr. \\
\midrule
2026-05-21 & 0 & 83 & -- & 9 & -- \\
2026-05-26 & 5 & 121 & 38 & 17 & 8 \\
2026-06-02 & 12 & 363 & 242 & 62 & 45 \\
2026-06-08 & 18 & 598 & 235 & 115 & 53 \\
2026-06-15 & 25 & 837 & 239 & 196 & 81 \\
2026-06-22 & 32 & 1{,}094 & 257 & 277 & 81 \\
2026-06-29 & 39 & 1{,}333 & 239 & 399 & 122 \\
2026-07-06 & 46 & 1{,}708 & 375 & 580 & 181 \\
2026-07-13 & 53 & 2{,}011 & 303 & 754 & 174 \\
2026-07-20 & 60 & 2{,}473 & 462 & 999 & 245 \\
2026-07-27 & 67 & 3{,}360 & 887$^\dagger$ & 1{,}487 & 488$^\dagger$ \\
2026-08-03 & 74 & 3{,}874 & 514 & 1{,}751 & 264 \\
2026-08-10 & 81 & 4{,}449 & 575 & 2{,}061 & 310 \\
2026-08-16 & 87 & 5{,}021 & 572 & 2{,}378 & 317 \\
2026-08-24 & 95 & 5{,}656 & 635 & 2{,}715 & 337 \\
2026-08-31 & 102 & 6{,}186 & 530 & 3{,}007 & 292 \\
2026-09-07 & 109 & 6{,}757 & 571 & 3{,}267 & 260 \\
2026-09-13 & 115 & 7{,}258 & 501 & 3{,}510 & 243 \\
2026-09-20 & 122 & 7{,}773 & 515 & 3{,}759 & 249 \\
2026-09-27 & 129 & 8{,}116 & 343 & 3{,}924 & 165 \\
2026-10-04 & 136 & 8{,}603 & 487 & 4{,}148 & 224 \\
2026-10-11 & 143 & 9{,}066 & 463 & 4{,}381 & 233 \\
\bottomrule
\end{tabular}

\vspace{2pt}
\par\footnotesize $^\dagger$The wk29$\to$wk30 increment contains a 22 July
retrospective reconciliation (272 cases, 174 deaths); the three intervals
13 July--3 August are merged into one 21-day interval in the likelihood.
\end{table}

The ECDC ``hospitalized in isolation'' series reads 896 (3 October), 945
(10 October), 899 (12 October) and 924 (14 October); the
residual $9{,}277-4{,}478-2{,}486=313$ (confirmed minus deaths minus
recorded recoveries) in the 14-October snapshot is far below the 924
reported in hospital, so the figure is read as an upper-bound proxy for
the hospitalized stock rather than a confirmed-case caliber. The weekly
series and the 14/16 October daily anchors are combined only in the
log-linear reference of Section~\ref{sec:results}, weekly points
entering as daily means and daily register intervals by length at their
midpoints.
\subsubsection{Data limitations}\label{sec:caliber}

Four features of the surveillance constrain the modeling and motivate
the observation layer of section~\ref{sec:obsmodel}. First, the
confirmed-case denominator is incomplete and revisable: the publicly
reported cumulative count fell from 8{,}728 (data to 6 October) to
8{,}665 (reported 7 October), and the report of 26 May listed 1{,}077
\emph{suspected} cases against 121 confirmed. A large share of deaths
occurs in the community (71\% of 372 in the period ending 8 October; about 60\% according to the WHO incident manager
), so the hospitalized share of cases is an estimated
parameter rather than a fixed input. Second, the increment series carries
two documented one-off registration transfers---the 22 July
retrospective reconciliation (272 cases, 174 deaths) inside the
wk29$\to$wk30 increment and the late-May suspected-to-confirmed
reclassification---which are absorbed by interval merging and by the
detection ramp respectively rather than modeled as noise. Third,
cross-source counts differ slightly (ECDC 4{,}410 deaths to 12 October
against UKHSA 4{,}381 to 11 October; 7{,}773/3{,}759 against
7{,}733/3{,}732 on 20 September) and are not spliced. Fourth, contact
follow-up is reported in incompatible calibers---flow (``about 45\%
followed up'' on 3 June, rising to 84.2\% by 12 August
) versus stock (ECDC's 73.1\% under
follow-up on 5 October)---which motivates the low
initial ring-coverage branch of the vaccination scenario.
\subsection{Compartmental model}\label{sec:model}

\subsubsection{State variables and equations}

The population is divided into susceptible $S$, incubating $E$,
symptomatic-community $I_C$ (patients who will not seek care),
care-seeking $I_S$ (patients progressing to hospitalization), hospitalized
$H$, recovered $R$, deceased $F$, and vaccinated $V$. Symptomatic
individuals are routed at symptom onset into the two tracks with
probabilities $p$ (care-seeking) and $1-p$ (community). The incubation
period is a five-stage linear chain---approximating the fitted gamma
latency of shape $\approx4.7$---with per-stage rate $c=5/7.68\approx0.651$
days$^{-1}$ (mean 7.68 days, SD 3.53 days, BDBV, $n=116$
\cite{Wamala2010}). The system is
\begin{subequations}\label{eq:model}
\begin{align}
\dot E_1 &= \lambda S/N_{\mathrm{eff}} - cE_1, \qquad
\dot E_i = c(E_{i-1}-E_i),\quad i=2,\dots,5,\\
\dot I_S &= p\,cE_5 - h I_S,\\
\dot I_C &= (1-p)\,cE_5 - (\mu_I+\gamma_I) I_C,\\
\dot H   &= h I_S - (\gamma_H+\mu_H) H,\\
\dot R   &= \gamma_I I_C + \gamma_H H,\\
\dot F   &= \mu_I I_C + \mu_H H,\\
\dot S   &= -\lambda S/N_{\mathrm{eff}},
\qquad\text{with }\lambda=\beta(t)\,(I_S+I_C+\kappa H).
\end{align}
\end{subequations}
where $h=0.198\,$d$^{-1}$ is the care-seeking hazard (onset-to-hospitalization
median 3.5 days \cite{VanKerkhove2015}; for an exponential sojourn this
implies a mean of 5.05 days, so $h=1/5.05$), $\mu_I+\gamma_I=0.095\,$d$^{-1}$
is the community exit rate (mean residence 10.5 days, from median
onset-to-death 9--10 days and onset-to-recovery 10 days
\cite{VanKerkhove2015} by moment matching under a gamma residence
distribution \cite{Lloyd2001,Wallinga2007}), split as $\mu_I=0.071$,
$\gamma_I=0.024$ to reproduce the community death share at an overall
case fatality $\bar c=0.60$, and $1/(\gamma_H+\mu_H)=8.84$ days is the mean
hospital residence time (survivor discharge 10.4 d and fatal course
7.0 d \cite{WHOEbolaResponse2014,Legrand2007}; cohort ranges move the
weighted mean only within $[8.3,9.7]$ d, and the Monte Carlo draws
$m_H\sim U[7,12]$ d span the wider variation). With the community-death
share $s_{\mathrm{comm}}=0.71$, the hospital and community case
fatalities are $c_1=(1-s_{\mathrm{comm}})\bar c/p=0.40$ and
$c_2=s_{\mathrm{comm}}\bar c/(1-p)=0.75$, so that
$p\,c_1+(1-p)\,c_2=\bar c=0.60$. Table~\ref{tab:priors} collects all
parameters with prior ranges and sources. The effective mixing
population is $N_{\mathrm{eff}}\in[8\times10^6,2.4\times10^7]$
(bottom-up: 65 affected health zones times an Ituri mean of
$1.94\times10^5$ inhabitants $\approx1.26\times10^7$
\cite{SSHAP2026,KhethImpilo2026}); the central value $1.35\times10^7$
adds the populous North Kivu driver zones to the Ituri core. The
transmission rate decays as $\beta(t)=\beta_0 e^{-\rho t}$ during
calibration, re-expressed as explicit interventions in the scenario
analysis (section~\ref{sec:scenarios}). A vaccinated class $V$
(section~\ref{sec:vacc}) routes a fraction $\nu(t)$ of the force of
infection on $S$ into $V$; vaccinated individuals transmit and die at
reduced rates (anti-infection efficacy $\mathrm{VE}_S$, anti-death
efficacy $\mathrm{VE}_D$), the class being omitted from \eqref{eq:model}
to keep the transmission analysis self-contained.
\subsubsection{Observation layer}\label{sec:obsmodel}

Confirmed cases are not infections. We model the cumulative confirmed
count as
\begin{equation}\label{eq:obs}
C(t) \;=\; d(t)\,\textstyle\int_0^t \lambda(\tau)\, S(\tau)/N_{\mathrm{eff}}\;d\tau, \qquad
d(t) = \frac{d_{\max}}{1+\exp\!\big(-(t-t_0)/s\big)},
\end{equation}
with $d_{\max}\in[0.30,0.95]$ the plateau detection fraction
(consistency domain $[0.34,0.95]$; Section~\ref{sec:results}); $d_{\max}=0.75$
is fixed exogenously from the under-reporting anchor, not estimated from the
case-count likelihood, and at $t=148$ the logistic ramp has effectively reached
this plateau so $d(148)\approx d_{\max}=0.75$. The ramp
absorbs the joint effect of case ascertainment, laboratory throughput
and surveillance scale-up, letting the observed early rate
$r_0\approx0.109\,$d$^{-1}$ coexist with a much slower intrinsic growth.
Weekly confirmed increments are modeled as negative-binomial realizations
with mean equal to the increment of $C(t)$; the dispersion is fixed at
$k_{\mathrm{obs}}=18.44$, estimated by the Pearson method of moments over
the 23 increments of the extended window from the relation
$\hat k=\sum_j\mu_j^2\big/\bigl[\sum_j(\Delta C_j^{\mathrm{obs}}-\mu_j)^2-\sum_j\mu_j\big]$
($\mu_j$ the increment of the fitted $C(t)$; calendar interval lengths
do not enter because the likelihood is over increments), iterated jointly
with the fitted mean structure to a fixed point---the post-climb window
$t\ge60$ alone gives $17.61$ and wk38--wk41 alone $29.45$---with
sensitivity over $\{20,39,18.4,\infty\}$.
\paragraph{Confirmation-timing convention and its robustness.}
The cumulative convention of \eqref{eq:obs} is equivalent to a nowcasting
correction: accrued infections are re-counted at the detection fraction
current at the report date, so the modeled series is directly comparable
with the retrospective cumulative snapshots the registers publish, and
the death-side offset ($+4.1\%$ between $d\cdot F(148)$ and the
registered deaths) is attributed to the death-register reporting lag.
A flow-caliber refit (confirmation at the time of infection) gives
$\beta_0=0.155\,$d$^{-1}$ ($-13\%$), $\Reff_0=1.50$, an earlier and
steeper ramp ($t_0=3.4$ d, $s=2.5$ d, detection saturating by
$t\approx32$ d) and a likelihood better by $2.3$ nats, the gain
concentrated in the earliest intervals; the detected-death offset grows
from $+4.1\%$ to $+9.6\%$ while the fatality reconciliation closes at
$d_{\max}\approx0.78$ against $0.81$ under the primary convention. We
therefore retain the cumulative (nowcasting) convention as primary and
report the flow-caliber refit as a sensitivity bound: the growth-phase
$\Reff_0$ spans $1.50$--$1.73$ across conventions, the model-free
current-transmission estimates being unaffected. The late-ramp
competitor of section~\ref{sec:ident} also disappears under the flow
convention (no basin beyond $t_0=40$ d within
$\Delta\mathrm{NLL}<5$), so the early-ramp reading does not hinge on
the confirmation-timing convention.
\subsubsection{The disorder-coupled variant}\label{sec:variant}

A variant considered in earlier work couples a social-disorder index $X$
to the isolation and transmission parameters. To test whether the data
require this coupling, we add the couplings with $\theta_a=\theta_b=0.4$
and an exogenously constructed $X(t)$, retaining the same free
parameters; Section~\ref{sec:results} reports the controlled comparison.

\subsection{Theoretical properties}\label{sec:theory}

\begin{proposition}[Positivity and boundedness]\label{prop:pos}
The closed positive octant is forward-invariant for system
\eqref{eq:model} with the vaccination inflow of section~\ref{sec:vacc},
and all solutions are bounded:
$S+\sum_{i=1}^{5}E_i+I_S+I_C+H+R+V\le N_{\mathrm{eff}}$ and $0\le F\le
\bar c\cdot N_{\mathrm{eff}}$, where $\bar c$ is the overall case
fatality.
\end{proposition}
\begin{proof}[Proof sketch]
\begin{sloppypar}
Each right-hand side is Lipschitz and points inward on the boundary
planes; the sum of all living-compartment derivatives equals
$-\lambda S/N_{\mathrm{eff}}-\mu_I I_C-\mu_H H\le0$, so the living total
is non-increasing (vaccination only moves mass between $S$ and $V$), and
$F$ is non-decreasing and bounded by $\bar c\,N_{\mathrm{eff}}$.
\end{sloppypar}
\end{proof}

\begin{proposition}[Next-generation reproduction number]\label{prop:ngm}
\begin{sloppypar}
At the disease-free equilibrium the next-generation matrix for the
infected classes $(E_1,\dots,E_5,I_S,I_C,H)$ has spectral radius
\end{sloppypar}
\begin{equation}\label{eq:R0}
\Reff_0=\beta_0\Bigg[\;\underbrace{\frac{p}{h}}_{\text{care-seeking track}}
+\underbrace{\frac{1-p}{\mu_I+\gamma_I}}_{\text{community track}}
+\underbrace{\frac{p\,\kappa}{\gamma_H+\mu_H}}_{\text{hospital pathway}}\;\Bigg].
\end{equation}
\end{proposition}
\begin{proof}[Proof sketch]
$F$ has nonzeros only in the $E_1$ row; $V^{-1}$ maps $E_1$ through the
five-stage chain (factor $1/c$ at the $I_S/I_C$ entry), then along the
care-seeking track (mean $1/h$ in $I_S$, every patient entering $H$ for a
mean $1/(\gamma_H+\mu_H)$ at reduced rate $\kappa$) or the community
track (mean $1/(\mu_I+\gamma_I)$); the spectral radius of a rank-one
$FV^{-1}$ is its unique nonzero eigenvalue.
\end{proof}

Two structural remarks follow. (i) In a single-$I$ formulation exited at
rate $q$, memorylessness forces all onset-to-exit delays to equal $1/q$;
with $q^{-1}=9.0$ days the model matches onset-to-death by accident while
overestimating onset-to-hospitalization by a factor $9.0/3.5\approx2.6$,
a deficiency the path-mixed structure removes by construction. (ii) With
the central parameters the bracket decomposes as $p/h=2.17$ ($22.4\%$),
$(1-p)/(\mu_I+\gamma_I)=6.00$ ($61.9\%$) and
$p\,\kappa/(\gamma_H+\mu_H)=1.52$ ($15.7\%$): the community track
dominates, and omitting the hospital pathway understates $\Reff_0$ by
$15.7\%$.

\paragraph{Care-seeking hazard correction.}
The care-seeking hazard was originally entered as $h=1/3.5=0.2857\,$d$^{-1}$,
treating the reported median onset-to-hospitalization delay of 3.5 days
as a mean; for an exponential sojourn the median implies a mean of
$3.5/\ln 2=5.05$ days, so $h=0.198\,$d$^{-1}$. Re-fitting separates the
correction from the window extension: $\beta_0$ falls from $0.1899$ to
$0.1784\,$d$^{-1}$ while the next-generation multiplier rises from
$9.026$ to $9.692$, so $\Reff_0$ moves from $1.714$ to
$1.729$---a net $+0.8\%$ revision, not the $+7.4\%$ a fixed-$\beta_0$
comparison would suggest (the displayed $1.729$ is the three-figure
rounding of the headline estimate $1.7289$, so this sub-per-cent shift is
the artifact of the median-to-mean hazard conversion rather than a change
in the transmission estimate). All results use $h=0.198$; the two hazard
calibers are indistinguishable in likelihood ($\Delta\mathrm{NLL}=0.013$ nats), so the correction is structural---anchored in the median-to-mean conversion---rather than data-driven, and the conclusion is unchanged.
\subsection{Calibration and model comparison}\label{sec:estimation}

Weekly confirmed increments over the calibration window are fitted by
maximum likelihood under the negative-binomial observation model of
section~\ref{sec:obsmodel}, with dispersion fixed at $k_{\mathrm{obs}}=18.44$
(the method-of-moments value over the full extended window, iterated to
a fixed point with the fitted means). The fits are
repeated over $k_{\mathrm{obs}}\in\{20,39,18.4,\infty\}$, with the AIC gap
between the two detection hypotheses $+45.5$, $+83.0$, $+42.1$ and
$+400.1$: the ramp hypothesis wins at every dispersion. Two hypotheses
are compared at fixed $p=0.43$: H1 (detection ramp; $\beta_0,\rho,d_{\max},t_0,s$)
and H2 (constant detection; $\beta_0,\rho,d$). Uncertainty is quantified
by profile likelihood over $(\beta_0,\rho)$, with $p$ and $(\mu_I,\gamma_I)$
varied along the case-fatality manifold, structural parameters drawn from
their prior ranges (Table~\ref{tab:priors}) and Monte Carlo propagation
($N=1000$); the disorder-coupled variant of section~\ref{sec:variant}
is compared against H1 at identical free parameters, and the data-alone
log-linear estimate ($\Reff\approx1.09$, 95\% CI $[1.04,1.14]$, derived
$T_g=12$ d) is reported as a model-free reference.
\begin{table}[htbp]
\centering\footnotesize\setlength{\tabcolsep}{4pt}
\caption{Parameters, values and prior ranges. Sources in the text.}\label{tab:priors}
\begin{tabular}{l>{\raggedright\arraybackslash}p{0.17\textwidth}>{\raggedright\arraybackslash}p{0.16\textwidth}>{\raggedright\arraybackslash}p{0.41\textwidth}}
\toprule
Symbol & Meaning & Central value & Range / source \\
\midrule
$c$ & incubation per-stage rate & $0.6510\,$d$^{-1}$ & mean 7.68 d (CI 7.04--8.32), SD 3.53, $n=116$ \cite{Wamala2010}\\
$h$ & care-seeking hazard & $0.198\,$d$^{-1}$ & median 3.5 d \cite{VanKerkhove2015}; exponential sojourn $\Rightarrow$ mean 5.05 d; sensitivity over Roddy IQR 1--8 d gives $h\in[0.125,1.0]$\\
$\mu_I$ & community death rate & $0.071\,$d$^{-1}$ & derived: community CFR $c_2=s_{\mathrm{comm}}\bar c/(1-p)=0.75$ \\
$\gamma_I$ & community recovery rate & $0.024\,$d$^{-1}$ & derived \\
$\gamma_H+\mu_H$ & hospital exit rate & $0.1131\,$d$^{-1}$ & weighted 8.84 d \\
$\mathrm{CFR}_H$ & hospital case fatality & 0.40 & $[0.35,0.50]$; $c_1=(1-s_{\mathrm{comm}})\bar c/p$ \\
$\kappa$ & hospital transmission factor & 0.40 & $[0.23,0.56]$ \cite{Nosocomial2026}\\
$p$ & care-seeking share & 0.43 (fixed) & $\{0.35,0.43,0.50\}$; derived from the community-death share \\
$T_g$ & generation time & 11.3 d (SD 5.1) & \emph{derived}: first-transmission convention $7.68+1/(\beta_0+\mu_I+\gamma_I)$ (SD from $\sqrt{3.5^2+3.7^2}=5.1$ d), reported at 12 d; model-implied size-biased value $17.5$ d (SD 10.5); sensitivity 10--18 d; EBOV serial interval 15.3 d (SD 9.1) \cite{WHOEbolaResponse2014} as contrast \\
$\Reff$ (2007) & published BDBV reference & 1.55 & 95\% CrI 1.02--2.36; pre-intervention CFR 31.4\% (21.8--42.8) \cite{DePadua2026} \\
$k_{\mathrm{obs}}$ & NB dispersion (counts) & 18.44 (fixed) & MoM (Pearson-type, 23 intervals, model means, iterated to a fixed point): $18.43$ over the full extended window, $17.61$ over the post-climb window $t\ge60$, $29.45$ over wk38--41 alone; sensitivity $\{20,39,18.4,\infty\}$: $\Delta$AIC $+45.5/+83.0/+42.1/+400.1$ \\
$N_{\mathrm{eff}}$ & effective mixing population & $1.35\times10^7$ & $[8,24]\times10^6$, bottom-up \cite{SSHAP2026,KhethImpilo2026}\\
$d_{\max}$ & plateau detection fraction & 0.75 (informed) & non-identifiable from case counts (Section~\ref{sec:results}); fixed from under-reporting $\approx15\%$ anchor \cite{DePadua2026}; fatality reconciliation closes at $d_{\max}\approx0.81$; varied over $[0.34,0.95]$ \\
\bottomrule
\end{tabular}
\end{table}

\subsubsection{Numerical implementation}\label{sec:numerics}

The system~\eqref{eq:model} is integrated with a fixed-step fourth-order
Runge--Kutta scheme (step 0.5 day; refinement to 0.1 day changes the
fitted $\beta_0$ by less than 0.1\%). Maximum-likelihood calibration uses
the Nelder--Mead simplex restarted from three initial simplexes, with
convergence declared when the simplex side length falls below $10^{-6}$;
the reported optimum is one of six basins located by a stratified
354-start global search (section~\ref{sec:ident}).
The initial condition is set by the early anchors of section~\ref{sec:data}
($\mathrm{cum\_inf}(0)=83/d(0)$, $d(0)\approx0.25$). Scenario projections
reuse the fitted $(\beta_0,\rho)$ and replace the decay with explicit
intervention schedules (section~\ref{sec:scenarios}): scenario B0 holds
transmission at its end-of-calibration value $\beta_0e^{-\rho\cdot148}\approx0.088$,
and scenario A extrapolates the fitted $\rho$. Monte Carlo intervals use
a per-sample refit: for each of $N=1000$ draws (seed convention
$20261011+7919i+13$) the model is re-optimized against the same
likelihood under the drawn structural parameters, the detection plateau
$d_{\max}\sim U[0.34,0.95]$ and the dispersion drawn equiprobably from
$\{12,18.44,39,\infty\}$ (this set includes the central MoM estimate
$k_{\mathrm{obs}}=18.44$ as one of its four equiprobable draws). All computations use
Python 3 with SciPy.
\subsubsection{Identifiability and confounding}\label{sec:ident}

Several identifiability results follow from fitting the negative-binomial
likelihood
\begin{equation}
\mathcal{L}(\boldsymbol{\theta})
= \prod_{j}\mathrm{NB}\!\big(\Delta C^{\mathrm{obs}}_{j};\,
\mu_{j}(\boldsymbol{\theta}),\,k_{\mathrm{obs}}\big),
\qquad
\mu_{j}=C(t_{j+1};\boldsymbol{\theta})-C(t_{j};\boldsymbol{\theta}),
\label{eq:ident-lik}
\end{equation}
to the reported confirmed-incidence increments, with dispersion fixed at
$k_{\mathrm{obs}}=18.44$: epidemic \emph{scale} and the \emph{level} of case
detection are statistically entangled, whereas the \emph{time-varying}
detection profile and the growth trajectory are separately recovered.

\paragraph{Structural non-identifiability of the detection ceiling
$d_{\max}$.}
The observation model links the true cumulative infection $Y(t)$ to the
reported count through a monotone, time-varying detection function:
\begin{equation}
C(t)=d(t)\,Y(t),\qquad
d(t)=\frac{d_{\max}}{1+\exp\!\big(-(t-t_{0})/s\big)}
=d_{\max}\,\sigma(t;t_0,s),
\label{eq:ident-obs}
\end{equation}
so that the reported increment depends on the product of the
detection profile and the infection inflow,
$d_{\max}\,\sigma(t)\,Y(t;\beta_{0},\rho,\ldots)$. Because only the
reported scale $C(t)$ is observed, $d_{\max}$ and $\beta_0$ act as
alternative multipliers of its height (lowering $d_{\max}$ by $a$ and
raising $\beta_0$ by $1/a$ preserves $C(t)$, $Y$ responding inversely to
$\beta_0$), so the likelihood is flat along the $(\beta_0,d_{\max})$
ridge: fixing $d_{\max}=0.75$ and letting it vary freely in
$[0.30,0.95]$ yield transmission rates differing by less than
$4\times10^{-5}$, the ceiling settling at the lower bound
($d_{\max}=0.30$) with the penalized likelihood unchanged (AIC gap
$+2.0$ for the extra parameter). The case-count likelihood thus carries no
information about $d_{\max}$, which we fix from an external
under-reporting anchor (approximately $15\%$ under-reporting implies
$d_{\max}\approx0.75$) and vary over $[0.30,0.95]$ in sensitivity
analysis.

\paragraph{What the case-count data fix.}
The early incidence trajectory is governed by $\Reff_0$ and the
generation interval, which depend on $\beta_0$ and the natural-history
parameters but not on how $\Reff_0$ factorizes. In the disorder-coupled
variant ($h_{\mathrm{eff}}=h(1-\theta X)$,
$\beta_{\mathrm{eff}}=\beta(1+\theta X)$), $\Reff_0=\beta_0\,G(X_0)$
with $G$ increasing in $\theta$; fitting $\theta\in\{0,0.2,0.4,0.6\}$
with the same five incidence parameters re-optimized, the fitted
$\beta_0$ decreases as $\theta$ increases so that $\beta_0\,G(X_0)$
changes little, while the likelihood improves monotonically. Even so,
the AIC gap between the uncoupled model M$_{A}$ ($\theta=0$) and the
weakly coupled M$_{B}$ ($\theta=0.2$) is below two units, so the
disorder coupling is statistically indistinguishable from zero:
$\Reff_0$ and $\theta$ are not separately identifiable from case-count
data alone. The same degeneracy appears between $\beta_0$ and the
waning rate $\rho$: a 200-replicate parametric bootstrap gives 95\%
intervals $[0.149,0.205]\,\mathrm{d}^{-1}$ and
$[0.0027,0.0063]\,\mathrm{d}^{-1}$ (profile intervals conditional on the
early-ramp basin are 32--39\% narrower and, for $t_0$, distorted by a
grid-endpoint artifact), so we report the bootstrap intervals as the
main caliber. Unlike the sparse-window fit, whose interval for $\rho$
straddled zero, the extended window excludes zero at both ends: the
growth-phase points sharpen the estimate of the waning rate rather than
reverse its sign, and the quoted $\Reff_0=1.729$ (bootstrap 95\%
$[1.44,1.98]$) rests on a positive, precisely estimated decay. The
bootstrap correlation $\mathrm{corr}(\beta_0,\rho)=+0.98$ implies that
$\Reff_0$ and $\Reff(148)$ are strongly negatively correlated ($-0.88$),
and the empirical width of the $\Reff(148)$ interval is one third of
what an independent combination of the two marginal intervals would
give, so the two intervals must not be read as independent evidence.
The detection-ramp midpoint $t_0$ is weakly constrained
($[1.0,14.9]$ d) whereas the steepness $s$ is better identified
($[4.3,13.7]$ d).

\paragraph{Global optimality and a competing basin.}
A stratified global search---354 starts (256 Latin-hypercube plus 98
directed) over the box $\beta_0\in[0.03,0.60]$,
$\rho\in[0,0.020]$, $t_0\in[-20,140]$, $s\in[3,60]$, each refined to
its basin bottom---locates six distinct basins, 217 of the starts
converging to the early-ramp optimum reported here; no start attains a
lower penalized likelihood, and an extended box (158 starts) reproduces
the same two leading basins. One competing basin is comparable in
likelihood: a late-ramp reading ($t_0=104.3$ d, $s=33.6$ d,
$\rho=0.067\,\mathrm{d}^{-1}$ outside the search box,
$\beta_0=0.443\,\mathrm{d}^{-1}$) costs only $\Delta\mathrm{NLL}=+1.92$, the
two basins giving nearly identical confirmed trajectories and
differing only in how the early growth is decomposed between
transmission and detection. We exclude it on four grounds: its
$\Reff_0=4.29$ exceeds the ceiling above which filovirus estimates are
not considered credible, its $\Reff(60)=0.08$ requires transmission to
be extinguished by late June although cases and deaths grew through
September with no intervention inside the window, it over-predicts the
mid- and late-window detected deaths by $15\%$ and $24\%$ (the
early-ramp reading under-predicts by $18\%$ and matches the tail), and
it requires $7{,}465$ true infections within the first twelve days
against $363$ reported ($95\%$ under-ascertainment). The same late-ramp
basin recurs under the uncorrected-hazard convention
($t_0=104.3\,\mathrm{d}$, $\Delta\mathrm{NLL}=+1.86$, $\Reff_0=4.14$) with
identical exclusion grounds, so the multi-modality is not an artefact
of the care-seeking correction. The detected-death increments, which lie outside the calibration
likelihood, cannot separate the two basins
($\Delta\mathrm{NLL}=-0.08$), so the exclusion rests on mechanistic
plausibility: this is the structural identifiability limit of a
single-logistic observation layer, flagged as a limitation
(section~\ref{sec:limits}). Under the adopted uniform-mixing,
single-logistic specification the data do not favor either the late-ramp
or the disorder-coupled alternative---the late-ramp reading is excluded
on mechanistic grounds (its likelihood penalty is only
$\Delta\mathrm{NLL}=+1.92$) and the disorder coupling is statistically
indistinguishable ($\Delta\mathrm{AIC}=-0.22$)---so the early-ramp,
uncoupled reading is retained by parsimony and mechanistic plausibility
rather than by likelihood alone. Fixing the ramp midpoint in the
intermediate range $t_0\in[40,90]$ finds no local minimum
($\Delta\mathrm{NLL}=+11.3$ at $t_0=50$), so a July--August
detection-ramp reading is rejected by the data; dropping the two
earliest intervals, which are contaminated by a suspected-to-confirmed
reclassification, leaves the early-ramp basin intact
($\Reff_0=1.70$, $-1.4\%$).

\paragraph{Identifiable structure.}
Although the detection \emph{ceiling} is unidentifiable, the \emph{shape}
is not: the logistic-ramp model H1$'$ against constant detection H2$'$
is penalized at two units per free parameter, so
$\Delta\mathrm{AIC}=+42.07$ reflects the likelihood
alone---$\sigma(t;t_0,s)$ modulates the bend location and steepness of
the reported increment trajectory, features not confounded with
$\beta_0$---and the improving-detection signal is a genuine feature of
the data. A 1000-sample Monte Carlo over the natural-history priors
yields one-year projection intervals dominated by scenario
assumptions---disorder persistence and the vaccination
specification---rather than by calibration uncertainty, while the daily
case anchors pin the absolute reporting level.
\section{Results}\label{sec:results}

\subsection{Model fit and comparison}

Table~\ref{tab:calib} lists the fitted parameters (central variant:
$p=0.43$ fixed, informed detection plateau $d_{\max}=0.75$). The
calibration window extends to $t=152$ (20 October), combining the
declaration anchor, 21 weekly points (wk21--wk41) and four daily
register intervals into 23 intervals. The fit reproduces the cumulative
confirmed count to $-2.4\%$ at $t=148$ (model 9{,}191 against observed
9,415), with a mean absolute interval deviation of 6.9\% in cases and
14.2\% in detected deaths; the detected-death caliber carries a baseline
offset of $+4.1\%$ ($d\cdot F(148)=4{,}721$ against 4,536), dominated by
the death-register reporting lag. The residuals are not white: the model
over-predicts the first interval ($+115\%$), under-predicts the
mid-epoch (mean $-8.5\%$ in cases and $-18.2\%$ in detected deaths over
$t=67$--$109$, the weeks following the 22 July reconciliation), and
converges to $-2.5\%/+0.4\%$ over the last four weeks; interval-level
diagnostics show no serial or dispersion structure (Durbin--Watson
$1.93$; runs test $p=0.57$; Ljung--Box $p=0.61$; mean Pearson $\chi^2$
$0.67$), so the pattern is a low-frequency mean mismatch that the
four-parameter ramp cannot absorb rather than a mis-specified
likelihood. The modeled symptomatic and hospitalized prevalence
($I_S+I_C+H=1{,}168$) exceeds the ECDC in-patient anchor of 924
($+26\%$), but the two are not the same caliber---the model counts all
symptomatic individuals, the report only hospitalized patients---and a
caliber-matched reading cuts both ways: the confirmed-caliber in-patient
count $d(148)\,H=284$ is 69\% below the report, while under the
low-detection reading ($d_{\max}=0.30$) the all-hospitalized census
rises to $H=948$---2.5 times the central $H=379$, as required by
$\mathrm{cum\_inf}(0)\propto 1/d_{\max}$, and 2.6\% above the
anchor---and the register arithmetic (924 in hospital against a
same-snapshot confirmed residual of 313) supports the second reading.
The two anchors are in tension---the hospital census favors $d_{\max}$
as low as $0.30$, the fatality reconciliation below closes at
$0.81$---so we treat $d_{\max}$ as an informed structural parameter,
report the central value $0.75$ over the domain $[0.34,0.95]$, keep the
weakly observed prevalence outside the likelihood, and flag the tension
as a limitation (section~\ref{sec:limits}). Figure~\ref{fig:weekly}
displays the fit.
\begin{table}[htbp]
\centering\footnotesize
\caption{Fitted parameters (central variant) and model comparison.
Point estimates are reported to three significant figures and are
reproducible from the deposited code. Intervals are 95\% parametric
bootstrap (200 replicates).}
\label{tab:calib}
\begin{tabularx}{\textwidth}{l>{\raggedright\arraybackslash}p{0.16\textwidth}>{\raggedright\arraybackslash}X}
\toprule
Parameter & Estimate & Notes \\
\midrule
$\beta_0$ & $0.1784\,$d$^{-1}$ & $0.149$--$0.205$ \\
$\rho$ & $0.00478\,$d$^{-1}$ & $0.0027$--$0.0063$ \\
$d_{\max}$ & 0.75 (informed) & non-identifiable from case counts; fixed from under-reporting anchor \\
$t_0$ (detection midpoint) & $5.57$ d & $1.0$--$14.9$ \\
$s$ (detection width) & $7.90$ d & $4.3$--$13.7$ \\
$\Reff_0$ (growth phase) & $1.729$ & $1.44$--$1.98$; nowcasting convention, $1.50$ under the flow convention (section~\ref{sec:obsmodel}); BDBV 2007: $1.55\,[1.02,2.36]$ \cite{DePadua2026} \\
\midrule
\multicolumn{3}{l}{AIC comparison (window to $t=152$; $k_{\mathrm{obs}}=18.44$):}\\
H1' (detection ramp, 5 param) & $130.11$ & \\
H2' (constant detection, 3 param) & $153.15$ & $\Delta$AIC $=+42.07$ \\
\multicolumn{3}{l}{Disorder coupling (same free parameters):}\\
M$_A$ (no coupling, $\theta=0$) & $130.11$ & \\
M$_B$ (coupled, $\theta=0.2$) & $130.00$ & $\Delta$AIC $=-0.22$, indistinguishable \\
\bottomrule
\end{tabularx}
\end{table}

Two findings follow. The detection-plateau hypothesis is strongly
favored ($\Delta\mathrm{AIC}=+42.07$ against constant detection): the
early exponential rate of confirmed cases ($r_0\approx0.109\,$d$^{-1}$)
is absorbed by the ramp from $d\approx0.25$ to the plateau rather than
by intrinsic transmission. The two quantities that remain free are
pinned by \emph{external} constraints: the care-seeking share $p$ is
unidentifiable from confirmed counts alone (the free fit drifts to
$p\to1$, contradicting the observed community share of deaths,
68--92\%), and the plateau detection fraction $d_{\max}$ carries
\emph{no} information from the case-count likelihood---with
$\mathrm{cum\_inf}(0)=83/d(0)$ the trajectory factorises as
$C(t)=83\,[\sigma(t)/\sigma(0)]\,G(t)$, in which $d_{\max}$ cancels
exactly, the infected subsystem being effectively linear over the
calibration window (susceptible depletion $\approx0.2\%$); the profile
likelihood over $d_{\max}\in[0.30,0.95]$ is flat to within $10^{-3}$ in
NLL (letting $d_{\max}$ vary freely settles at the lower bound
$d_{\max}=0.30$ and moves the re-fitted $\beta_0$ by less than
$4\times10^{-5}$), and $\Reff_0=1.729$ does not drift with the
detection ceiling. A reconciliation of the fatality accounts with the
register sets the consistency domain: the structural eventual CFR from
the compartment split,
$0.43\times0.40+0.57\times0.75=0.60$, exceeds the registered
$4{,}536/9{,}415=0.482$ because cases are under-ascertained by
$d_{\max}$ and deaths accrue with delay (only 65\% of eventual deaths
recorded by mid-October; full-caliber deaths grow from $6{,}295$ to
$9{,}705$); combining the corrections gives eventual
$\mathrm{CFR}\approx0.482\,d_{\max}/0.6486$, which is $0.557$ at
$d_{\max}=0.75$ and closes at $d_{\max}\approx0.81$ against the
structural $0.60$, while the lag-free form ($0.482\,d_{\max}$) would
have excluded $d_{\max}<0.45$ but with the lag $d_{\max}=0.50$ implies
$0.37$, inside the historical BDBV band $[0.22,0.43]$. These
eventual-CFR figures follow definition~(i): ``eventual deaths'' here
include mortality from infections arising after $t=148$ under scenario~A
run to convergence, so the $0.557$ at $d_{\max}=0.75$ and the
$[0.22,0.43]$ band are both definition-i quantities; the cohort
definition~(ii), which confines the denominator to the infections already
present by day~148, gives the separate values $0.854$ and $0.423$ and is
not mixed with the figures above. The reported
domain $d_{\max}\in[0.34,0.95]$ (the Monte Carlo sampling range) rests
on the identifiability grounds; it implies cumulative infections of
$9{,}911$--$27{,}691$ by 16 October (central $12{,}553$). The modeled
detected-death caliber over-states the registered death count by $+4.1\%$
at $t=148$ ($d\cdot F(148)=4{,}721$ against $4{,}536$); the
eventual-CFR correction above inherits this conservative $\sim+4\%$
caliber inconsistency, biasing the estimate upward rather than downward.
Two external anchors pull the ceiling downward---the hospital census above, and the
WHO Emergency Committee's assessment that the true number of infections
may be three to four times the reported count, implying
$d_{\max}\approx0.25$--$0.33$---so the central value $0.75$ sits
between the poles: it is fixed from the external under-ascertainment
evidence (roughly 15\% for the 2007 Uganda outbreak \cite{DePadua2026},
and the fatality closure at $0.81$), and the scale-dependent outputs are
reported across the full domain.
\begin{figure}[htbp]
\centering
\includegraphics[width=0.85\textwidth]{fig_v6_weekly_fit.png}
\caption{Fit of the observation-layer model to the verified
confirmed-case series. (a) Interval confirmed increments (log scale):
the declaration anchor (circle), UKHSA weekly points (squares), the
merged 21-day reconciliation interval (diamonds), the raw
reconciliation-week increment (cross) and daily register intervals
(triangles) against the fitted model (line). (b) Cumulative confirmed
cases: model (line) against the verified anchors. Central variant
$p=0.43$, $d_{\max}=0.75$ informed, $k_{\mathrm{obs}}=18.44$.}
\label{fig:weekly}
\end{figure}

\subsection{Current transmission state}

The fitted trajectory implies a declining effective reproduction number
$\Reff(t)=\beta_0e^{-\rho t}\times9.69$ \cite{Cori2013}, from $1.73$ in
the growth phase to $0.85$ at $t=148$; the crossing of one falls at day
115 (parametric-bootstrap 95\% interval $[106,137]$). The growth-phase
data also sharpen the decay rate: $\rho$ moves from $0.0019$ to
$0.0048\,\mathrm{d}^{-1}$ and its interval from $[-0.001,0.005]$ to
$[0.003,0.006]$---a precision gain rather than a reversal, since the
earlier point estimate already sat at the upper edge of the earlier
interval. The model-free log-linear reference is the appropriate
independent check, but it must exclude the growth phase, where the
detection ramp dominates the reported slope: over the full window the
interval increments give $r=+0.0070\,\mathrm{d}^{-1}$ and
$\Reff(12\,\mathrm{d})=1.09$, whereas the post-reconciliation windows
give $r=-0.0035$ ($t\ge67$) and $r=-0.0049$ ($t\ge60$), hence
$\Reff(12\,\mathrm{d})=0.96$ (95\% interval $[0.93,0.99]$) and $0.94$
($[0.91,0.98]$). The intervals rest on 14--16 adjacent, autocorrelated
points, so they are indicative rather than precise confidence
statements; the two-point extrapolation (wk39$\to$wk41, 343$\to$463)
yields an apparent $\Reff\approx1.29$ that the full series does not
support. The generation time has no BDBV-specific estimate and the two
derived conventions differ by more than the estimation error: the
full-sojourn value $17.5$ d (SD $10.5$) against the conservative
first-transmission value $T_g=11.3$ d (SD $5.1$), with the EBOV serial
interval ($15.3$ d, SD $9.1$ \cite{WHOEbolaResponse2014}) between them.
The point estimate stays within $[0.92,0.97]$ across the
post-reconciliation windows over $T_g\in[10,16]$ d, and every
post-reconciliation interval excludes one: the data now separate
decline from stability, though not decline from slow decline. Figure~\ref{fig:reff} summarizes the estimates. 
\begin{figure}[htbp]
\centering
\includegraphics[width=0.85\textwidth]{fig_v6_reff.png}
\caption{Effective reproduction number. The curve is the model estimate
$\Reff(t)=\beta_0e^{-\rho t}\times9.69$ (v6 fit, E1), marked at
$\Reff_0=1.73$ and $\Reff(148)=0.85$ with the unit crossing near day
115. The horizontal references are model-free log-linear estimates at
$T_g=12$ d under alternative window choices: the full calibration
window ($1.08$, contaminated by the detection ramp) and the
post-reconciliation windows $t\ge60$ ($0.94$), $t\ge67$ ($0.96$) and
$t\ge74$ ($0.95$).}
\label{fig:reff}
\end{figure}

As an out-of-sample check, the model calibrated only to the 21 weekly
points (to 11 October) and extrapolated five days gives 9{,}317
cumulative confirmed cases against 9,415 observed ($-1.0\%$) and 4{,}893
detected deaths against 4,536 ($+7.9\%$, reporting-lag dominated), at
52.5 cases/day against 69.8 observed ($-24.8\%$). The direction is
mechanical: the four low-leverage daily-tail intervals pull the
full-window fit toward a later, weaker ramp, lowering the fitted waning
rate $\rho$ (removing them raises $\beta_0$ by 10\% and $\rho$ by 29\%)
and thereby steepening the modeled late decline, so the weekly-only
calibration under-predicts the observed end-of-window daily rate by a
quarter. The five-day extrapolation nonetheless lands within 1\% on the
cumulative count---a miss that the single-day-versus-interval caliber of
the register partly explains and that we report rather than tune away.
\subsection{One-year projections}\label{sec:scenarios}

Table~\ref{tab:proj} juxtaposes the model projection with the
data-driven band that uses no model (daily rates of 69 confirmed cases
and 31 deaths, extrapolated under decline, zero growth and the fitted
post-reconciliation log-linear trend). Under the detected-death caliber
$d\cdot F$, scenario B0 (transmission held at its end-of-calibration
value) adds $+4{,}785$ deaths over one year and $+6{,}684$ confirmed
cases. Both lie below the data-driven band ($+5{,}200$ to $+11{,}315$
deaths; $+11{,}575$ to $+25{,}185$ cases, the upper edge set by the
zero-growth assumption): the fitted decay of $\beta$
steepens the model's decline relative to every data-driven rate
assumption, the decline row included, so the two envelopes need not
coincide and the model projection is reported for comparison, not as a
forecast. The band spans a factor of 2.2; anchoring the fitted trend at
$t\ge60$ rather than $t\ge67$ gives $+11{,}673$ confirmed cases and
$+5{,}244$ deaths, close to the decline row. Figure~\ref{fig:scen} places the
model trajectories against the weekly and daily-tail anchors.

\emph{Uncertainty caveat.} The Monte Carlo intervals redraw the
natural-history parameters and $d_{\max}\sim U[0.34,0.95]$ with the
corrected care-seeking hazard $h=0.198$ (section~\ref{sec:theory}),
redraw the negative-binomial dispersion equiprobably from
$\{12,18.44,39,\infty\}$ (the central MoM estimate included as one of the
four draws), and re-optimize the calibrated parameters per
draw ($N=1000$); they cover calibration and structural uncertainty but
not the low-frequency mean mismatch of section~\ref{sec:results} or the
observation-layer form. Because the full-death projection is the most
$d_{\max}$-sensitive calibrated quantity, the tabulated $d\cdot F$
intervals are the conservative summary (the full-death $F$ caliber's B0
97.5\% bound reaches roughly twice the median).
\begin{figure}[htbp]
\centering
\includegraphics[width=0.85\textwidth]{fig_v6_scen.png}
\caption{One-year projections from $t=148$ (2026-10-16). Model
scenarios A (transmission continues to decay with the fitted $\rho$),
B0 (transmission held at its end-of-calibration value
$\beta_0e^{-\rho\cdot148}\approx0.088$) and the disorder-coupled B,
which coincides with B0 in the central uncoupled variant, against the
weekly and daily-tail anchors. (a) Cumulative confirmed cases (log
scale). (b) Cumulative deaths in the detected $d\cdot F$ caliber; the
projection is anchored to the observed register by the end-of-window
completeness $4{,}649/d\cdot F(152)=0.95$ (a constant reporting
offset), so that it continues from the registered-death series, while
the increments tabulated in Table~\ref{tab:proj} are in the unscaled
$d\cdot F$ caliber.}
\label{fig:scen}
\end{figure}

\begin{table}[htbp]
\centering
\caption{One-year increments (horizon $t=513$, 365 days after
$t=148$=2026-10-16). Model rows use the central variant ($p=0.43$,
$d_{\max}=0.75$, uncoupled $\theta_a=\theta_b=0$) calibrated through
20 October; deaths in the detected
$d\cdot F$ caliber, comparable with the registered confirmed-death
series. Monte Carlo intervals ($N=1000$) redraw the natural-history
parameters, $d_{\max}\sim U[0.34,0.95]$ and $k_{\mathrm{obs}}$, and
re-optimize the calibrated parameters per draw. Data-driven rows use no
model (daily rates of 69 confirmed cases and 31 deaths).}
\label{tab:proj}
\begin{tabularx}{\textwidth}{Xrr}
\toprule
Assumption & Confirmed cases & Deaths \\
\midrule
Model: transmission continues to decay (A) & $+2{,}920$ & $+2{,}558$ \\
\quad MC 95\% interval & $[+2{,}047,\,+2{,}960]$ & $[+1{,}876,\,+2{,}824]$ \\
Model: transmission held at end-of-calibration value (B0) & $+6{,}684$ & $+4{,}785$ \\
\quad MC 95\% interval & $[+4{,}107,\,+6{,}965]$ & $[+3{,}030,\,+5{,}332]$ \\
Data-driven: decline $r=-0.005$/d & $+11{,}575$ & $+5{,}200$ \\
Data-driven: zero growth & $+25{,}185$ & $+11{,}315$ \\
Data-driven: fitted trend (log-linear, $t\ge67$) & $+14{,}159$ & $+6{,}361$ \\
\bottomrule
\end{tabularx}
\vspace{2pt}
\par\small Model death rows in the full-death $F$ caliber (all deaths
counted, not only detected ones) are $+3{,}410$ (A) and $+6{,}380$
(B0), with Monte Carlo medians $+3{,}791$ $[2{,}285,7{,}263]$ and
$+6{,}844$ $[3{,}681,13{,}514]$---the full-death medians sit above
their deterministic counterparts because the draws include $d_{\max}$
values below the central 0.75, which rescale the infection total by
$1/d_{\max}$. Monte Carlo medians of the detected-caliber model rows
are $+2{,}884$ confirmed and $+2{,}510$ deaths (A), and $+6{,}671$ and
$+4{,}709$ (B0). Both calibers and the coupled scenario B---which
coincides numerically with B0 in the uncoupled central variant, because
$X$ does not enter its dynamics---are reported in the reproducibility
repository.
\end{table}

\subsection{Vaccination scenario}\label{sec:vacc}

No vaccine is licensed for the Bundibugyo species, so the
ring-vaccination scenario demonstrates the framework's ability to
evaluate a hypothetical rVSV-ZEBOV-like campaign rather than a deployed
countermeasure; the reductions reported are illustrative lower bounds.
The efficacy pair, $\mathrm{VE}_S=0.15$ (susceptibility) and
$\mathrm{VE}_D=0.40$ (death given infection), is a demonstration
assumption for a Zaire-targeted vaccine applied to Bundibugyo virus, for
which neither cross-protection nor ring effectiveness has been observed.
Against the baseline B0 with the occupational campaign that has reached
$8{,}623$ frontline workers---$7{,}775$ under a protocol-governed
compassionate-use program granted after the DRC's request to the
International Coordinating Group  and $848$ under
the subsequently suspended BRAVO protocol---by 7
October 2026 (the demonstration initializes $V(0)=9{,}000$ as a rounded
stand-in), a prospective ring program ($c_{\mathrm{ring}}=20$ contacts
traced per confirmed case, $f_{\mathrm{ring}}=0.5$) reduces the deaths
added over the projection year (horizon $t=513$) by only $0.20\%$, the
same value in the full-death $F$ and detected-death $d\cdot F$ calibers
(the two coincide because $d(t)$ is at its plateau over the projection
year). The full grid over
$f_{\mathrm{ring}}\in\{0.2,0.5,1.0\}$ and
$c_{\mathrm{ring}}\in\{10,20,30\}$ spans $0.048$--$0.566\%$ (central
$0.202\%$; Monte Carlo median $0.147\%$, 95\% interval $[0.060,0.303]$).
This small effect must not be
read as evidence that vaccination is ineffective: the occupational
campaign protects frontline workers and preserves the health-system
capacity on which isolation and safe burial depend---a contribution the
compartmental structure cannot represent---while the ring program is
evaluated purely on its direct effect under uniform mixing, where the
binding constraint is $f_{\mathrm{ring}}$, the unobserved fraction of
transmission inside traceable rings, so the modeled reduction is an
illustrative lower bound rather than a statement about efficacy. The
reason is structural: ring dosing is proportional to daily confirmed
incidence (of order 60 cases/day), and the simulated campaign reaches
only $114{,}000$ vaccinated individuals by day 513---$0.85\%$ of
$N_{\mathrm{eff}}$, the occupational cohort included---too small a
coverage to move
population-level transmission under uniform mixing. A quantitatively
meaningful vaccination scenario requires a far higher dosing rate or an
explicit high-exposure subpopulation; because anti-death efficacy
dominates the outcome, cumulative deaths remain the primary endpoint.
Figure~\ref{fig:vacc} contrasts the vaccinated and unvaccinated one-year
death counts under the central variant.
\begin{figure}[htbp]
\centering
\includegraphics[width=0.85\textwidth]{fig_v6_vacc.png}
\caption{One-year cumulative deaths in the detected-death caliber under
the central ring-vaccination variant (D; $c_{\mathrm{ring}}=20$,
$f_{\mathrm{ring}}=0.5$) against the unvaccinated baseline B0
(transmission held at its end-of-calibration value): $9{,}670$ versus
$9{,}680$ deaths at $t=513$, i.e. a $0.20\%$ reduction in the deaths
added over the projection year ($4{,}775$ versus $4{,}785$), the same
reduction as in the full-death caliber. The effect stays below $1\%$
under uniform mixing, because anti-death efficacy dominates the
outcome.}
\label{fig:vacc}
\end{figure}

\subsection{Comparison with previous Ebola outbreaks}\label{sec:compare}

The 2026 DRC Bundibugyo epidemic is the first BDBV outbreak to reach
cross-border scale; Table~\ref{tab:compare} places it against the
recorded epidemics that define the field's scale and methods---the two
prior BDBV outbreaks and the two largest of any species---asking how
much of an outbreak's burden is set by intrinsic transmissibility and
how much by context. All historical figures are published reports,
independent of the model outputs.
\begin{table}[htbp]
\centering
\caption{Comparison of the 2026 DRC Bundibugyo outbreak with previous
Ebola epidemics. Reproduction numbers are pre-intervention or growth-phase
effective estimates as reported by the cited sources; CFR is among
reported cases; ``--'' indicates no directly comparable published
estimate. Response summarizes the dominant control measure available at the
time.}
\label{tab:compare}
\resizebox{\textwidth}{!}{%
\begin{tabular}{llllllll}
\toprule
Outbreak & Country & Species & Cases & Deaths & CFR & $\Reff$ & Response \\
\midrule
Uganda 2007 & Uganda & BDBV & 149 & 37 & 25--34\% & 1.55 $[1.02,2.36]$ & isolation \\
DRC 2012 & DRC & BDBV & 57 & 29 & $\approx$51\% & -- & isolation \\
West Africa 2014--16 & GIN/LBR/SLE & EBOV & 28{,}646 & 11{,}323 & 39.5\% & 1.5--2.5 & global NPI + ETU \\
DRC Kivu 2018--20 & DRC & EBOV & 3{,}470 & 2{,}287 & $\approx$66\% & -- & ring vax (rVSV) \\
DRC 2026 (this study) & DRC & BDBV & 9{,}668 & 4{,}649 & 48.2\% & 1.73 $[1.44,1.98]$ & none licensed\footnotemark \\
\bottomrule
\end{tabular}}
\end{table}
\footnotetext{No vaccine is licensed for Bundibugyo virus. The
vaccination scenario (D) in this study is therefore a methodological
demonstration of the model's counterfactual capability rather than a
description of a deployed intervention; its assumed efficacy is an
illustrative upper bound, not an established estimate.}

The same-species benchmark is Uganda 2007, the only other BDBV epidemic
with a published control-adjusted estimate: $\Reff=1.55$ (95\% CrI
$1.02$--$2.36$) falling to $0.50$ after control measures, with an
adjusted pre-intervention CFR of $31.4\%$ ($21.8$--$42.8$) and roughly
$15\%$ under-ascertainment \cite{DePadua2026}. Our growth-phase
estimate $\Reff_0=1.73$ lies inside that interval, so the difference in
scale is not a difference in the virus: by 16 October 2026 the outbreak
had recorded $9{,}415$ confirmed cases and $4{,}536$ deaths, about $63$
times the Uganda case count and $123$ times its deaths, despite a
statistically compatible estimated $\Reff_0$ (1.73 against 1.55). Uganda 2007 was contained within
months in a single district with an intact surveillance and isolation
system \cite{Wamala2010}; the 2026 outbreak spans eight provinces and 66
health zones in an active conflict zone with large-scale displacement
\cite{SSHAP2026}. The higher recorded CFR ($48.2\%$ versus Uganda's
crude $25$--$34\%$) is consistent with more infections never reaching a
treatment centre rather than with a more virulent virus.

West Africa 2014--2016 frames the upper bound of what Ebola can do:
country-specific reproduction numbers of $1.5$--$2.5$
\cite{Chowell2004,Nishiura2014} that rose rather than fell as the virus
entered densely populated capitals, delayed detection, super-spreading
(only $20\%$ of cases generated at least $73\%$ of onward infections
\cite{Kucharski2015}), and the most severe under-ascertainment of any
Ebola epidemic, which left cumulative curves unusable until
observation-layer and nowcasting corrections were applied
\cite{Pandey2014,Merler2015}. The present outbreak's $\Reff_0=1.73$
sits inside that range, yet West Africa still produced
about $2.5$ times our recorded deaths: scale is set by response speed
and detection rather than $\Reff$ alone. The closest precedent in
setting is Kivu 2018--2020: the same country, an active conflict zone,
community distrust and attacks on response staff, strongly overdispersed
transmission ($k=0.29$; $17.8\%$ of cases generating $80\%$ of
transmission) \cite{Soubrier2026}, and $1{,}171$ survivors. The decisive
difference is that Kivu had a \emph{species-matched} ring-vaccination
campaign---more than $303{,}000$ doses of rVSV-ZEBOV-GP with novel
therapeutics and community engagement, to which the decline is widely
attributed \cite{WHOKivu2020}; our vaccination analysis
(section~\ref{sec:vacc}) finds the mirror image with a Zaire-targeted
vaccine against Bundibugyo virus.

Taken together, the comparators isolate what governs scale: the 2026
outbreak pairs transmissibility statistically compatible with Uganda
2007 ($\Reff_0=1.73$ inside $[1.02,2.36]$) with the
West-Africa/Kivu-level context, which is why its burden is large despite
an unexceptional $\Reff$ and its projection is an interval spanning a
factor of 2.2---West Africa shows that a higher $\Reff$ plus delayed
detection yields continental spread, Kivu that conflict alone does not
preclude control once a matched vaccine and community engagement
coincide, and Uganda 2007 that the same virus in a contained setting
stays small, while by 16 October its $4{,}536$ deaths already exceeded
the entire Kivu epidemic ($2{,}287$). Each comparator also returns to
the methodological point our model formalizes: under-ascertainment must
be corrected before reproduction numbers can be read off cumulative
curves, and naive fits that treat confirmed cases as infections
overstate $\Reff$ by a factor $1.54$ ($2.66$ against $1.73$), a generic
property of Ebola surveillance rather than a feature of 2026.
\section{Discussion}\label{sec:discussion}

The methodological contribution reaches past the 2026 outbreak. Fitting
\emph{cumulative} curves invites a reproduction-number estimator
structurally bounded below by one
($\hat{\Reff}=\exp[T_g\ln(C_2/C_1)/\Delta t]\ge1$): a zero-growth
incidence profile already produces $\hat{\Reff}\approx1.09$--$1.14$ at
this outbreak's parameters ($C=9{,}415$ at the calibration endpoint
with daily incidence $69.8$/d gives $i/C=0.0074\,$d$^{-1}$, hence
$\hat{\Reff}=\exp(12\times0.0074)=1.09$ at $T_g=12$ d and
$\exp(17.5\times0.0074)=1.14$ at $T_g=17.5$ d), so cumulative-ratio
statements of the form ``still above one'' are uninformative here, and
incidence-based estimation with an interval is the minimum standard.
The early acceleration of confirmed cases is joint evidence about
transmission and detection: a model without an observation layer
attributes all of it to transmission, and the fitted consequence is
quantitative---the constant-detection baseline H2$'$ overstates the
intrinsic $\Reff_0$ by a factor $1.54$ ($2.66$ against $1.73$,
$\Delta\mathrm{AIC}=+42.07$). The path-mixed symptomatic structure is
likewise a correctness requirement, not a refinement: it is the minimal
modification that lets a compartmental model hold both the care-seeking
delay and the community residence time that the natural history fixes.

The behavioral coupling question deserves an explicit burden of proof.
At identical free parameters, the model without the coupling fits the
verified weekly series as well as the disorder-coupled variant
($\Delta\mathrm{AIC}=-0.22$); this does not show that social disorder is
irrelevant---the $X(t)$ construction is exogenous and hand-weighted, and
under the strongly coupled variant ($\theta=0.4$) the disorder
amplification of a stagnant scenario reaches a factor $1.47$ in
cumulative cases by day 513 ($1.28$ at $\theta=0.2$)---because $X(t)$ is
constructed to decay over the projection horizon, the amplification
converges to a finite factor rather than growing without bound, and the
converged end-of-horizon multiples ($1.47$ and $1.28$) sit in a
$1.3$--$1.5\times$ range that tracks the coupling strength---but it does show
that the present data do
not \emph{require} the coupling, and any claim that it drives the
epidemic must be supported by a fit improvement that survives the
comparison.

Two external anchors frame the fitted values: the only published
BDBV-specific estimate, from Uganda 2007, is $\Reff_t=1.55$ with a
pre-intervention CFR of $31.4\%$
\cite{DePadua2026}, which our growth-phase estimate ($1.73$) lies
inside; and the Kivu transmission-chain reconstruction finds strong
individual overdispersion ($k=0.29$, $17.8\%$ of cases generating
$80\%$ of transmission) \cite{Soubrier2026}, implying that a
deterministic model overestimates the probability of large final sizes
and that stochastic fade-out is a live alternative no deterministic fit
can exclude.
\subsection{Limitations}\label{sec:limits}

(i) The register is incomplete and non-monotone: suspected-to-confirmed
ratios near nine in May, a suspected-to-confirmed reclassification in
late May and a retrospective reconciliation on 22 July, both absorbed
into the negative-binomial dispersion rather than modeled explicitly
(the reconciliation window is merged, and the reclassification intervals
excluded, in sensitivity analyses). The verified weekly series now spans
wk21--wk41 plus four daily register intervals, which constrains the
growth phase directly, but the earliest intervals remain the least
reliable and the mid-epoch increments are under-predicted by $8.5\%$ in
cases and $18\%$ in detected deaths---a low-frequency mismatch the
four-parameter ramp cannot absorb.
(ii) The single-population formulation cannot represent the observed
spatial heterogeneity (48 of 64 affected health zones with recent cases
and 16 extinguished); the projection is a national
aggregate, biased in a known direction---pooling extinguished and active
zones dilutes the local reproduction number, so the national estimate is
likely conservative for the zones that still drive the epidemic. (iii)
The observation layer carries informed assumptions rather than
identified quantities: the detection ceiling $d_{\max}$ is fixed from an
under-reporting anchor while the external anchors pull apart---the
hospital census and the WHO IHR Emergency Committee's three- to
four-fold under-ascertainment statement  point to
$0.25$--$0.33$, whereas the fatality reconciliation closes at $0.81$---
with the central value $0.75$ set between them; the
confirmation-timing convention is an assumption, the growth-phase
$\Reff_0$ spanning $1.50$--$1.73$ across the two conventions
(section~\ref{sec:obsmodel}); and the
care-seeking share $p$ and the community death share are linked through
the same case-fatality decomposition, so the sensitivity grid over them
probes the feasibility of the symptomatic split rather than independent
uncertainty. (iv) The vaccine scenario relies on cross-protection
assumptions that are plausible but unproven (an illustrative upper
bound, with a zero-default cross-protection multiplier). (v) The model
has no explicit funeral (post-death) transmission pathway: unsafe
burials sustained a substantial share of transmission in prior
outbreaks, and the community-death share of 68--92\% here falls into
that unmodeled channel, so the recommendation to cut community
transmission is anchored in practice to safe burial and community
engagement rather than to a mechanism in the equations. (vi)
Transmission-level overdispersion is represented only in the
observation-layer negative binomial, not in the infection process: given
the strong individual overdispersion documented for Kivu, a
deterministic model overstates the probability of large final sizes and
stochastic fade-out remains a live alternative. Asymptomatic or mild
Bundibugyo infections are not modeled, so if a meaningful fraction of
infections are sub-clinical the symptomatic-only $\Reff_0$ is a lower
bound and the cross-species comparison is not strictly like-for-like.
(vii) The analysis is descriptive rather than operational and its equity
dimension is not quantified: it does not translate the transmission
estimates into decision-trigger logic a response team could act on, and
it treats the effective population uniformly, so it cannot locate where
vaccination or isolation capacity is most needed. (viii) The likelihood
surface is multi-modal: a late-ramp reading of the same single-logistic
observation layer attains $\Delta\mathrm{NLL}=+1.92$ and is excluded on
mechanistic grounds rather than by the data (section~\ref{sec:ident}),
and a two-component detection function improves the AIC by $4.5$ while
implying a third growth narrative ($\Reff_0=0.70$); the reported
intervals are conditional on the adopted early-ramp basin and
single-logistic form and therefore understate the structural uncertainty
in $\Reff_0$. (ix) The out-of-sample check is honest about its miss:
calibrated to the weekly points alone, the model reproduces the
cumulative count to $1\%$ but under-predicts the end-of-window daily
incidence by a quarter ($52.5$ against $69.8$ cases/day), a
single-day-versus-interval caliber difference of the register that we
report rather than tune away.
\subsection{Conclusions}

The 2026 DRC Bundibugyo outbreak is the largest epidemic ever caused by
this virus species, yet its intrinsic transmissibility is unexceptional:
the growth-phase $\Reff_0=1.73$ (bootstrap 95\% $[1.44,1.98]$) lies
inside the credible interval of the only other quantified BDBV outbreak
(Uganda 2007) and at the lower end of the Zaire-ebolavirus range. Its
cross-border scale is explained by context---prolonged transmission in a
conflict-affected, mobile population with a high community-death share
and only partial case ascertainment---rather than by a more
transmissible pathogen, and the path-mixed SEIHRD structure with an
explicit detection ramp provides a transparent, reproducible baseline
that separates the two. The current effective reproduction number is
estimated below one ($0.85$, bootstrap 95\% $[0.77,0.97]$; the
model-free post-reconciliation estimate $0.96$ $[0.93,0.99]$ agrees).
The growth-phase and current-reproduction bootstrap intervals are strongly
negatively correlated ($\mathrm{corr}\approx-0.88$, inherited from the
$\beta_0$--$\rho$ correlation) and are not independent estimates; see
Section~\ref{sec:ident}.
The one-year projection band spans a factor of $2.2$; the response
should therefore be judged on its ability to hold the decline while
raising detection and cutting community transmission, above all through
safe and dignified burial---the mechanism the model cannot itself
represent (Limitation (v)) but the one through which the epidemic is
most likely to be interrupted.
\subsection{Code and data availability}\label{sec:availability}

The analysis is fully reproducible from deposited code and the public
surveillance series used in this study. The model and the calibration, Monte-Carlo
propagation and scenario pipeline are implemented in Python 3
(\texttt{numpy}, \texttt{scipy}, \texttt{matplotlib}) and deposited in a
public version-controlled repository under an open license; the
repository URL and its Zenodo archived-snapshot DOI \zenododoi{} are
available at submission (the deposition was completed and archived with
Zenodo before submission). All inputs are public surveillance series published by WHO, the UK Health Security Agency (UKHSA), the European Centre for Disease Prevention and Control (ECDC), Africa CDC, the DRC Ministry of Public Health and associated news agencies, transcribed and archived by the deposited scripts; no individual-level or unpublished data were used, and every plotted or tabulated number is produced directly from these series. The full enumerated source list with URLs is provided in the deposited repository and as Supplementary Material.
\begin{thebibliography}{99}

\bibitem{SSHAP2026} Villa J, Wilkinson A, Bedford J, Schmidt-Sane M,
Duclos D. Contextual Note on the Ebola Bundibugyo Outbreak in Ituri
(2026). Social Science in Humanitarian Action Platform (SSHAP), 2026
(Ituri population about 7 million including roughly 1 million internally
displaced persons; conflict, mining-linked mobility and funeral
practices).
\url{https://www.socialscienceinaction.org/resources/contextual-note-on-the-ebola-bundibugyo-outbreak-in-ituri-2026}

\bibitem{KhethImpilo2026} Kheth'Impilo-DRC. Emergency Ebola outbreak
response, Congo DR: project report, 1 July--31 August 2026 (B\'eni health
zone, North Kivu: estimated population 631{,}268, area 21{,}677
km$^2$).
\url{https://cl.globalgiving.org/projects/emergency-ebola-outbreak-response-congo-dr/reports}

\bibitem{Legrand2007} Legrand J, Grais RF, Boelle PY, Valleron AJ,
Flahault A. Understanding the dynamics of Ebola epidemics. \emph{Epidemiol
Infect}. 2007;135(4):610--621.

\bibitem{Cori2013} Cori A, Ferguson NM, Fraser C, Cauchemez S. A new
framework and software to estimate time-varying reproduction numbers
during epidemics. \emph{Am J Epidemiol}. 2013;178(9):1505--1512.

\bibitem{VanKerkhove2015} Van Kerkhove MD, Bento AI, Mills HL, Ferguson
NM, Donnelly CA. A review of epidemiological parameters from Ebola
outbreaks to inform epidemiological modelling and microforecasting.
\emph{Sci Data}. 2015;2:150019. doi:10.1038/sdata.2015.19 (Table~5:
hospitalization median 3.5 d, range 0--8, $n=26$; onset-to-death medians
9--10 d; onset-to-recovery median 10 d).

\bibitem{Wamala2010} Wamala JF, Lukwago L, Malimbo M, et al. Ebola
hemorrhagic fever associated with novel virus strain, Uganda,
2007--2008. \emph{Emerg Infect Dis}. 2010;16(7):1087--1092.
doi:10.3201/eid1607.091525 (PMID 20587179; incubation-period summary
mean 7.68 days, 95\% CI 7.04--8.32, SD 3.53, $n=116$, transcribed via
\cite{Quilty2026}).

\bibitem{Quilty2026} Quilty BJ, et al. Estimating the effectiveness of
syndromic screening at airports for Bundibugyo ebolavirus disease.
medRxiv 2026.06.11.26355442v2 (preprint; transcribes the Wamala et al.\
2010 incubation-period summary, mean 7.68 days, SD 3.53, $n=116$).
doi:10.64898/2026.06.11.26355442

\bibitem{DePadua2026} de Padua B, Akhmetzhanov AR. Estimated
transmissibility and case fatality rate of Bundibugyo virus, Uganda,
2007. \emph{Emerg Infect Dis}. 2026;32(10):1762--1764.
doi:10.3201/eid3210.261175 (Research Letter: $\Reff_t=1.55$, 95\% CrI
1.02--2.36; pre-intervention CFR 31.4\%, 95\% CrI 21.8--42.8).
\url{https://wwwnc.cdc.gov/eid/article/32/10/26-1175}

\bibitem{Soubrier2026} Soubrier H, Seo D, Barks P, Flasche S, Camacho A,
et al. Quantifying the heterogeneity and determinants of Ebola Za\"ire
transmission during the 10th outbreak in DRC, 2018--2020. medRxiv
2026, doi:10.64898/2026.07.14.26358072 (preprint; $k=0.29$, 95\% CI
0.26--0.32; 17.8\% of cases generated 80\% of transmission; incidence-rate
ratios $4.34$ ($3.47$--$5.31$) for transmission from individuals who died
outside treatment centres and $1.79$ ($1.54$--$2.06$) for cases not
isolated in a treatment centre; EBOV Zaire----cross-species extrapolation
noted).

\bibitem{Nosocomial2026} Hospital transmission reduction factor
$\kappa\in[0.23,0.56]$: bracketed by triangulating the
transmission-outside-isolation incidence-rate ratios reported for the
2018--2020 DRC Kivu outbreak by Soubrier et al.\ \cite{Soubrier2026}
(IRR $4.34$ for deaths outside treatment centres, $1.79$ for cases not
isolated in a treatment centre), with the central value $\kappa=0.40$
chosen so that the hospital pathway contributes about one sixth of the
reproduction number (Section~\ref{sec:theory}).

\bibitem{WHOEbolaResponse2014} WHO Ebola Response Team. Ebola virus
disease in West Africa---the first 9 months of the epidemic and forward
projections. \emph{N Engl J Med}. 2014;371:1481--1495.
doi:10.1056/NEJMoa1411100 (serial interval 15.3 d, SD 9.1).

\bibitem{WHOKivu2020} WHO. 10th Ebola virus disease outbreak in the
Democratic Republic of the Congo declared over. 25 June 2020.
\url{https://www.who.int/news/item/25-06-2020-10th-ebola-outbreak-in-the-democratic-republic-of-the-congo-declared-over}

\bibitem{Lloyd2001} Lloyd AL. Realistic distributions of infectious
periods in epidemic models: changing patterns of persistence and
dynamics. \emph{Theor Popul Biol}. 2001;60(1):59--71.
doi:10.1006/tpbi.2001.1525.

\bibitem{Wallinga2007} Wallinga J, Lipsitch M. How generation intervals
shape the relationship between growth rates and reproductive numbers.
\emph{Proc R Soc B}. 2007;274(1609):599--604. doi:10.1098/rspb.2006.3754

\bibitem{Chowell2004} Chowell G, Hengartner NW, Castillo-Chavez C,
Fenimore PW, Hyman JM. The basic reproductive number of Ebola and the
effects of public health measures: the cases of Congo and Uganda.
\emph{J Theor Biol}. 2004;229(1):119--126. doi:10.1016/j.jtbi.2004.03.006.

\bibitem{Nishiura2014} Nishiura H, Chowell G. Early transmission dynamics
of Ebola virus disease (EVD), West Africa, March to August 2014.
\emph{Euro Surveill}. 2014;19(36):20894.
doi:10.2807/1560-7917.ES.2014.19.36.20894.

\bibitem{Pandey2014} Pandey A, Atkins KE, Medlock J, Wenzel N, Townsend
JP, Childs JE, Nyenswah TG, Ndeffo-Mbah ML, Galvani AP. Strategies for
containing Ebola in West Africa. \emph{Science}. 2014;346(6212):991--995.
doi:10.1126/science.1260612.

\bibitem{Kucharski2015} International Ebola Response Team. Exposure
patterns driving Ebola transmission in West Africa: a retrospective
observational study. \emph{PLoS Med}. 2016;13(11):e1002170.
doi:10.1371/journal.pmed.1002170.

\bibitem{Merler2015} Merler S, Ajelli M, Fumanelli L, Gomes MFC, Piontti
AP, Rossi L, Chao DL, Longini IM Jr, Halloran ME, Vespignani A.
Spatiotemporal spread of the 2014 outbreak of Ebola virus disease in
Liberia and the effectiveness of non-pharmaceutical interventions: a
computational modelling analysis. \emph{Lancet Infect Dis}.
2015;15(2):204--211. doi:10.1016/S1473-3099(14)71074-6.

\bibitem{Meltzer2014} Meltzer MI, Atkins CY, Santibanez S, Knust B,
Petersen BW, Ervin ED, Nichol ST, Damon IK, Washington ML. Estimating
the future number of cases in the Ebola epidemic---Liberia and Sierra
Leone, 2014--2015. \emph{MMWR Surveill Summ}. 2014;63(Suppl 3):1--14.

\end{thebibliography}

\end{document}
